% !TeX root = ../../Book5.tex
% 纲要标题：### B.3 表示论中的范畴化
% 纲要位置：工作记录/计划纲要.md，第 4593 行。
% #### B.3.1 Grothendieck 群与两种加法
% #### B.3.2 从张量恒等式到函子
% #### B.3.3 分次与复形中的符号
\section{表示论中的范畴化}
\label{b5:sec:B03}

一个表示环只记住同构类型的加法和乘法。回到表示、态射与函子以后，原来的整数恒等式可能来自更具体的直和分解或正合列。范畴化研究的正是这种提升；首先要说明取的是哪一种 Grothendieck 群，以及哪些信息在取群时被忘掉。

\subsection{Grothendieck 群与两种加法}
\label{b5:subsec:B03:01}

设 \(\mathcal A\) 为对象同构类构成集合的加性范畴。由对象符号 \([X]\) 生成，并加入
\([X\oplus Y]=[X]+[Y]\)
的交换群，称为分裂 Grothendieck 群。若 \(\mathcal A\) 还是交换范畴，则可以进一步对所有短正合列加入
\([M]=[L]+[N]\)，得到这里记为 \(G_0(\mathcal A)\) 的群。
\symidx{\(G_0(\mathcal A)\)：按短正合列关系定义的 Grothendieck 群}

\ProseParagraphBreak
两种关系不能混同。设 \(k\) 为特征 \(p\) 的域，考虑循环群 \(C_p\) 在 \(k\) 上的有限维表示。由 \(kC_p\cong k[t]/(t^p)\)，不可分解模为 \(M_d=k[t]/(t^d)\)、\(1\le d\le p\)。它们在分裂 Grothendieck 群中各给独立的基元，因为 Krull--Schmidt 分解唯一；而在 \(G_0\) 中，它们都满足
\([M_d]=d[M_1]\)，故只剩一个基元。这里 \(G_0\) 只记录简单合成因子的重数，分裂 Grothendieck 群仍可区分非分裂的不可分解对象。
\begin{proposition}\label{b5:prop:B-tensor-grothendieck}
设 \(\mathcal A\) 为对象同构类构成集合的有限长度交换范畴，配双加性且在每个变量上正合的张量积 \(\otimes:\mathcal A\times\mathcal A\to\mathcal A\)。记 \(G_0(\mathcal A)\) 为对对象同构类加入全部短正合列关系的 Grothendieck 群，\([X]\) 为对象 \(X\) 的类。则对 \(X,Y\in\mathcal A\)，乘法
\([X][Y]=[X\otimes Y]\)
在 \(G_0(\mathcal A)\) 上良定义；若有结合约束及单位对象，就得到相应的结合含幺环。其加法群以简单对象同构类为基。
\end{proposition}
\begin{proof}
固定一个对象作张量保持短正合列，所以对象乘法尊重定义群时加入的关系，双加性便使乘法下降。结合与单位同构给出环的结合律及单位。每个对象沿合成列等于其简单因子之和；Jordan--Hölder 保证重数不依赖列，所以取各简单重数给出到自由交换群的逆映射。
\end{proof}
Hopf 代数的有限维模范畴中，张量在底层域上正合，便给出这种环。这里的“有限长度”不意味着只有有限多个简单类型，也不意味着整个范畴半单。基础定义及例子可对照\cite[Definition 1.5.8; Section 4.5]{EtingofEtAl2015TensorCategories}。

\subsection{从张量恒等式到函子}
\label{b5:subsec:B03:02}

取 \(\mathfrak{sl}_2\) 的有限维复表示范畴，全部对象完全可约。以 \(L_m\)、\(m\ge0\) 为简单对象，则
\[
[L_1][L_m]=[L_{m+1}]+[L_{m-1}]
\quad(m\ge1),\qquad [L_1][L_0]=[L_1].
\]
这是\cref{b5:thm:sl2-clebsch-gordan}在复数域上取一个张量因子为 \(L_1\) 所得的分解，在 Grothendieck 环中的像。令 \(t=[L_1]\)，便可递推地把每个 \([L_m]\) 写成 \(t\) 的一个首一 \(m\) 次整数多项式。不同简单类线性无关，所以这些首一递推多项式也保证
\[
G_0\bigl(\operatorname{Rep}_{\mathrm{fd}}\mathfrak{sl}_2\bigr)
\cong\mathbb Z[t].
\]
在范畴中，乘以 \(t\) 被提升为函子
\(F(V)=L_1\otimes V\)，其在简单对象上的作用就是上述直和分解。对象以外还有交织映射、自然变换以及张量结合约束；这些信息在只取环时看不见。因此给出一个环同构，并没有自动构造范畴等价；给出一串正整数，也没有自动指定相应的自然变换及其关系。
\ProseParagraphBreak
例如 \([L_2]=t^2-1\) 在对象层面来自 \(L_1\otimes L_1\simeq L_2\oplus L_0\)，在函子层面则给出
\[
F^2\simeq (L_2\otimes-)\oplus\operatorname{Id}.
\]
具体地，翻转 \(\tau\) 与 \(\mathfrak{sl}_2\) 作用交换，投影 \((1+\tau)/2\)、\((1-\tau)/2\) 分别取出 \(\operatorname{Sym}^2L_1=L_2\) 与 \(\bigwedge^2L_1=L_0\)。将这两个投影张量任意对象的恒等映射，便得到与每个表示同态相容的自然投影。这是实际函子的直和分解；环中出现的减号只是在直和关系取 Grothendieck 类后移项得到的。

\subsection{分次与复形中的符号}
\label{b5:subsec:B03:03}

设 \(k\) 为域。有限维 \(\mathbb Z\)-分次 \(k\)-向量空间的分裂 Grothendieck 环可用 Laurent 变量 \(v\) 记录次数平移：
\[
[V]=\sum_{j\in\mathbb Z}(\dim V_j)v^j.
\]
张量积按次数相加，这个对应给出与 \(\mathbb Z[v,v^{-1}]\) 的环同构。每个分次片的维数都非负，而形式差允许负系数。

\ProseParagraphBreak
取有界的分次复形 \(C^\bullet\)，各 \(C^i\) 都有限维，且微分 \(d^i:C^i\to C^{i+1}\) 的内部次数为零，即 \(d^i(C^i_j)\subseteq C^{i+1}_j\)。于是圈 \(Z^i=\ker d^i\)、边界 \(B^i=\operatorname{im}d^{i-1}\) 和上同调 \(H^i(C^\bullet)\) 都继承内部的 \(\mathbb Z\)-分次，并有分次短正合列
\[
\begin{aligned}
0&\longrightarrow Z^i\longrightarrow C^i\longrightarrow B^{i+1}\longrightarrow0,\\
0&\longrightarrow B^i\longrightarrow Z^i\longrightarrow H^i(C^\bullet)\longrightarrow0.
\end{aligned}
\]
这两列逐内部次数分裂，故在上述分裂 Grothendieck 群中有
\([C^i]=[B^i]+[H^i(C^\bullet)]+[B^{i+1}]\)。乘以 \((-1)^i\) 后求和，边界项成对抵消，得到 \(\mathbb Z[v,v^{-1}]\) 中的 Euler 类恒等式
\[
\sum_i(-1)^i[C^i]=\sum_i(-1)^i[H^i(C^\bullet)].
\]

\ProseParagraphBreak
于是有两种性质不同的符号：内部次数平移产生 \(v\)，复形移位产生负号。把二者混为一谈，会把分次维数和 Euler 特征写错。本节用张量函子和分次空间介绍范畴化的基本想法；更深的范畴化还要求生成函子及其自然变换满足量子群或 Hecke 关系。
